% ===========================================================================
%  chapter2.tex  --  Глава 2. Экспериментальное исследование
% ===========================================================================

\newpage
\section{Экспериментальное исследование}
\label{sec:experiments}

В~данной главе описывается экспериментальный контур, разработанный для
сравнительного анализа классических и~нейросетевых моделей прогнозирования
временных рядов. Последовательно изложены: постановка эксперимента
(данные, модели, метрики), результаты прямого сравнения на~данных M3 и~M4
Competition, анализ гибридного подхода Helformer, эксперименты с~итеративным
прогнозированием и~предобучением, а~также обобщающее обсуждение
полученных результатов.

% ---------------------------------------------------------------------------
\subsection{Постановка эксперимента}
\label{subsec:setup}

\subsubsection{Данные}
\label{sssec:data}

В~качестве экспериментальной базы используются два наиболее авторитетных
набора данных для одномерного прогнозирования: M3 Competition~\cite{makridakis2000m3}
и~M4 Competition~\cite{makridakis2020m4}.

\textbf{M3 Competition} содержит 3003~временных ряда,
разбитых на шесть частотных категорий. В~настоящей работе используются
три из~них: \textit{yearly} (645 рядов, горизонт прогноза $H{=}6$),
\textit{quarterly} (756 рядов, $H{=}8$) и~\textit{monthly}
(1428 рядов, $H{=}18$). Длина рядов варьируется от~14~наблюдений
для годовых до~126 для месячных, что создает различные режимы для
нейросетевых моделей: от~экстремально малых выборок до~умеренных.

\textbf{M4 Competition} содержит свыше 100\,000~рядов шести частот.
Из них в~эксперименте задействованы три частотные категории:
\textit{quarterly} (24\,000 рядов, $H{=}8$), \textit{monthly}
(48\,000 рядов, $H{=}18$) и~\textit{daily} (4227 рядов, $H{=}14$).
Включение ежедневных рядов позволяет исследовать поведение моделей
в~режиме длинных обучающих выборок (до~9919 наблюдений).

\textbf{Выборка.} Из~каждой категории случайным образом (seed=42)
отбирается по~30~рядов, что дает 90~рядов для M3 и~90~рядов для M4.
Такой объем обусловлен вычислительным бюджетом: обучение каждой
нейросетевой модели выполняется отдельно на~каждом ряде (per-series training),
и~полный прогон десяти моделей на~90~рядах занимает порядка 4--8~часов
процессорного времени. При этом 30~рядов на~категорию позволяют получить
устойчивые средние оценки метрик, хотя охват составляет лишь около
3--5\,\% от полных наборов.

\textbf{Разбиение.} Для каждого ряда последние $H$~значений
выделяются как тестовая выборка; все предшествующие наблюдения
используются для обучения. Этот протокол соответствует стандарту
соревнований M3/M4~\cite{makridakis2000m3, makridakis2020m4}
и~обеспечивает прямую сопоставимость с~опубликованными результатами.

\subsubsection{Модели и~их конфигурация}
\label{sssec:models}

В~эксперименте участвуют 11~моделей, которые можно разбить на~четыре
группы (табл.~\ref{tab:model-overview}).

\begin{table}[H]
\centering
\caption{Обзор моделей эксперимента}
\label{tab:model-overview}
\small
\begin{tabular}{@{}llll@{}}
\toprule
Модель & Тип & Декомпозиция & Ключевые параметры \\
\midrule
Seasonal Naive & baseline & --- & последний сезон \\
Auto-ARIMA     & классическая & разности & автоподбор $(p,d,q)$ \\
ETS            & классическая & фиксированная & автоподбор \\
Prophet        & классическая & фиксированная & additive/multiplicative \\
\midrule
DLinear        & нейросетевая & фиксированная (MA) & $L{=}512$, lr=0.001 \\
N-BEATS        & нейросетевая & обучаемая (trend+season) & стеки $\times$4, lr=0.001 \\
TimesNet       & нейросетевая & обучаемая (2D) & top-$k$=3, $d$=32, lr=0.001 \\
\midrule
Autoformer     & трансформерная & обучаемая (series decomp.) & $d$=64, $n_h$=4, lr=0.0001 \\
FEDformer      & трансформерная & обучаемая (MOEDecomp) & $d$=64, $n_h$=4, lr=0.0001 \\
PatchTST       & трансформерная & без декомпозиции & patch=8, $d$=64, lr=0.0001 \\
\midrule
Helformer      & гибридная & фиксированная (HW) & $d$=32, $n_h$=2, lr=0.001 \\
\bottomrule
\end{tabular}
\end{table}

\textbf{Режим обучения.} Все нейросетевые модели обучаются в~режиме
per-series: для каждого временного ряда модель инициализируется заново
и~обучается только на~данных этого ряда. Такой режим отличается от~режима
глобального обучения, при котором одна модель обучается сразу на~множестве
рядов~\cite{bandara2020forecasting}, и~заведомо ставит нейросетевые
модели в~невыгодное положение на~коротких рядах, однако позволяет
оценить способность архитектуры извлекать полезные паттерны из~ограниченных
данных --- ключевой вопрос настоящего исследования.

Количество эпох обучения составляет 80 для прямого прогнозирования
основных моделей и~40 для Helformer. Для ежедневных рядов M4
Autoformer и~FEDformer обучаются 30~эпох ввиду вычислительных ограничений.
Оптимизатор --- Adam~\cite{kingma2015adam} (lr=0.001 для DLinear, N-BEATS,
TimesNet, Helformer; lr=0.0001 для трансформерных моделей).

\textbf{Нормализация.} Перед подачей в~нейросетевые модели ряд
нормализуется по~$z$-score (вычитание среднего, деление на~стандартное
отклонение); для Helformer применяется Min-Max нормализация.
Обратное преобразование выполняется после генерации прогноза.

\subsubsection{Протокол оценки}
\label{sssec:protocol}

Эксперимент включает три протокола прогнозирования:

\begin{enumerate}[nosep]
\item \textbf{Прямой прогноз (direct):} модель обучается, после чего
  генерирует все $H$~шагов прогноза за~один проход. Это основной
  режим для всех сравнений.
\item \textbf{Итеративный прогноз (rollout):} модель прогнозирует
  один шаг вперед, прогнозное значение добавляется во~входное окно,
  и~процесс повторяется $H$~раз. Данный режим моделирует реальный
  сценарий эксплуатации.
\item \textbf{Предобучение (pretrain):} модель предварительно обучается
  на~всех рядах данной категории, затем дообучается на~конкретном ряде.
  Данный режим проверяет гипотезу о~недостатке данных при per-series
  обучении.
\end{enumerate}

\textbf{Метрики.} Для оценки качества прогноза используются шесть метрик:
\begin{itemize}[nosep]
\item sMAPE (symmetric Mean Absolute Percentage Error) --- основная метрика
  соревнований M3/M4~\cite{hyndman2006accuracy};
\item MAPE (Mean Absolute Percentage Error);
\item RMSE (Root Mean Squared Error);
\item MSE (Mean Squared Error);
\item MAE (Mean Absolute Error);
\item MASE (Mean Absolute Scaled Error)~\cite{hyndman2006accuracy} ---
  масштабно-инвариантная метрика, нормирующая ошибку на~ошибку наивного
  прогноза.
\end{itemize}

В~таблицах и~обсуждении основной акцент делается на~sMAPE
(процентная, безразмерная, допускает агрегацию по~рядам с~разным
масштабом) и~MASE (учитывает базовый уровень сложности ряда).

% ---------------------------------------------------------------------------
\subsection{Прямое сравнение моделей}
\label{subsec:direct}

\subsubsection{Результаты на~M3 Competition}
\label{sssec:m3-results}

Таблица~\ref{tab:m3-smape} содержит средние значения sMAPE по~категориям
и~в~целом для всех десяти моделей на~90~рядах M3. Жирным выделены лучшие
результаты в~каждой категории.

\begin{table}[H]
\centering
\caption{sMAPE (\%) на~M3 Competition: средние по~категориям (30 рядов на~категорию)}
\label{tab:m3-smape}
\small
\begin{tabular}{@{}lcccc@{}}
\toprule
Модель & Yearly & Quarterly & Monthly & Overall \\
\midrule
Seasonal Naive & 16.42 & 8.31 & 20.40 & 15.04 \\
Auto-ARIMA     & \textbf{12.64} & 8.36 & 26.38 & 15.79 \\
ETS            & 12.84 & \textbf{7.80} & 21.04 & \textbf{13.89} \\
Prophet        & 14.96 & 17.15 & 24.13 & 18.75 \\
\midrule
DLinear        & 37.64 & 19.32 & 24.89 & 27.28 \\
N-BEATS        & 27.37 & 11.58 & \textbf{20.58} & 19.84 \\
TimesNet       & 18.31 & 8.18 & 27.12 & 17.87 \\
\midrule
Autoformer     & 29.31 & 14.63 & 24.66 & 22.87 \\
FEDformer      & 27.75 & 16.91 & 26.64 & 23.77 \\
PatchTST       & 30.02 & 11.29 & 23.72 & 21.68 \\
\bottomrule
\end{tabular}
\end{table}

Таблица~\ref{tab:m3-mase} дополняет картину метрикой MASE,
нечувствительной к~масштабу ряда.

\begin{table}[H]
\centering
\caption{MASE на~M3 Competition: средние по~категориям}
\label{tab:m3-mase}
\small
\begin{tabular}{@{}lcccc@{}}
\toprule
Модель & Yearly & Quarterly & Monthly & Overall \\
\midrule
Seasonal Naive & 1.00 & 1.00 & 1.00 & 1.00 \\
Auto-ARIMA     & --- & --- & --- & 1.59 \\
ETS            & --- & --- & --- & \textbf{1.47} \\
Prophet        & --- & --- & --- & 1.86 \\
\midrule
DLinear        & --- & --- & --- & 4.24 \\
N-BEATS        & --- & --- & --- & 2.69 \\
TimesNet       & --- & --- & --- & 1.99 \\
\midrule
Autoformer     & --- & --- & --- & 3.06 \\
FEDformer      & --- & --- & --- & 3.14 \\
PatchTST       & --- & --- & --- & 2.78 \\
\bottomrule
\end{tabular}
\end{table}

Из~полученных данных следует ряд наблюдений.

\textit{Доминирование ETS.}
ETS уверенно занимает первое место по~общему sMAPE (13.89\,\%) и~MASE (1.47).
Этот результат согласуется с~выводами соревнований M3/M4, где
ETS традиционно входит в~число лучших одномодельных подходов~\cite{makridakis2000m3,
hyndman2008expsmoothing}.

\textit{Зависимость от~частоты.}
На~годовых рядах (самые короткие, 14--47~точек) нейросетевые модели
демонстрируют наибольшее отставание: даже лучший из~них TimesNet (18.31)
проигрывает Auto-ARIMA (12.64) более чем на~5~п.п. На~квартальных рядах
TimesNet (8.18) практически не~уступает ETS (7.80), а~PatchTST (11.29)
и~N-BEATS (11.58) находятся в~зоне умеренной конкурентоспособности.
На~месячных рядах N-BEATS (20.58) показывает лучший результат среди
всех моделей, опережая ETS (21.04), что свидетельствует о~способности
обучаемой декомпозиции N-BEATS эффективно работать при достаточном
объеме данных (более~100~наблюдений).

\textit{DLinear.}
Несмотря на~свою архитектурную простоту, DLinear показывает слабые
результаты на~M3 (27.28 overall), что объясняется малым объемом
обучающих данных: модель не~может эффективно разделить тренд
и~сезонность при per-series обучении на~коротких рядах.

\begin{figure}[H]
\centering
\includegraphics[width=\textwidth]{bar_m3_by_freq.png}
\caption{sMAPE по~категориям M3 Competition для всех моделей}
\label{fig:bar-m3-freq}
\end{figure}

На~рис.~\ref{fig:bar-m3-freq} визуально представлено распределение
sMAPE по~категориям. Четко видна закономерность: по~мере увеличения
длины рядов (yearly~$\to$~quarterly~$\to$~monthly) разрыв между
классическими и~нейросетевыми моделями сокращается.

\subsubsection{Результаты на~M4 Competition}
\label{sssec:m4-results}

Таблица~\ref{tab:m4qm-smape} содержит результаты на~квартальных
и~месячных рядах M4.

\begin{table}[H]
\centering
\caption{sMAPE (\%) на~M4 Competition: квартальные и~месячные ряды}
\label{tab:m4qm-smape}
\small
\begin{tabular}{@{}lccc@{}}
\toprule
Модель & Quarterly & Monthly & Overall \\
\midrule
Seasonal Naive & 17.39 & 18.70 & 18.05 \\
Auto-ARIMA     & 15.39 & 18.99 & 17.19 \\
ETS            & 16.72 & \textbf{15.40} & \textbf{16.06} \\
Prophet        & 22.03 & 20.46 & 21.24 \\
\midrule
DLinear        & 25.70 & 29.51 & 27.61 \\
N-BEATS        & 28.83 & 26.31 & 27.57 \\
TimesNet       & \textbf{18.17} & 28.03 & 23.10 \\
\midrule
Autoformer     & 34.17 & 33.86 & 34.02 \\
FEDformer      & 37.00 & 29.94 & 33.47 \\
PatchTST       & 25.00 & 28.87 & 26.94 \\
\bottomrule
\end{tabular}
\end{table}

На~M4 наблюдается более выраженное превосходство классических моделей.
ETS сохраняет лидерство (16.06), а~разрыв с~лучшей нейросетевой моделью
(TimesNet, 23.10) увеличивается до~7~п.п. Autoformer и~FEDformer
демонстрируют наихудшие результаты (34.02 и~33.47 соответственно),
что указывает на~неспособность обучаемой декомпозиции трансформерного
типа эффективно работать при ограниченных данных per-series обучения
на~рядах M4.

Принципиально иная картина наблюдается на~ежедневных рядах M4
(табл.~\ref{tab:m4d-smape}).

\begin{table}[H]
\centering
\caption{sMAPE (\%) на~M4 Competition: ежедневные ряды (30 рядов)}
\label{tab:m4d-smape}
\small
\begin{tabular}{@{}lcc@{}}
\toprule
Модель & sMAPE & MASE \\
\midrule
Seasonal Naive & 2.80 & 1.19 \\
Auto-ARIMA     & 2.45 & 1.02 \\
ETS            & 2.46 & 1.02 \\
Prophet        & 10.86 & 3.72 \\
\midrule
DLinear        & \textbf{2.43} & \textbf{1.00} \\
N-BEATS        & 3.36 & 1.39 \\
TimesNet       & 2.87 & 1.14 \\
\midrule
Autoformer     & 4.02 & 1.75 \\
FEDformer      & 3.43 & 1.55 \\
PatchTST       & 2.85 & 1.20 \\
\bottomrule
\end{tabular}
\end{table}

\begin{figure}[H]
\centering
\includegraphics[width=\textwidth]{bar_m4_daily_full.png}
\caption{sMAPE на~ежедневных рядах M4 для всех моделей}
\label{fig:bar-m4-daily}
\end{figure}

На~ежедневных рядах DLinear (2.43) впервые опережает ETS (2.46) и~Auto-ARIMA
(2.45), хотя различия минимальны. TimesNet (2.87) и~PatchTST (2.85)
показывают результаты, сопоставимые с~классическими моделями. Этот
результат демонстрирует, что при достаточной длине обучающей выборки
(сотни и~тысячи наблюдений) нейросетевые модели становятся полностью
конкурентоспособными с~классическими подходами. Autoformer и~FEDformer
обучались с~уменьшенным числом эпох (30 вместо 80), что частично
объясняет их более слабые результаты.

\subsubsection{Кросс-датасетный анализ}
\label{sssec:cross-dataset}

Для обобщения результатов табл.~\ref{tab:cross-smape} сводит средние
sMAPE по~всем шести экспериментальным категориям.

\begin{table}[H]
\centering
\caption{Кросс-датасетный sMAPE (\%): 10 моделей $\times$ 6 категорий}
\label{tab:cross-smape}
\small
\begin{tabular}{@{}lcccccc|c@{}}
\toprule
Модель & M3-Y & M3-Q & M3-M & M4-Q & M4-M & M4-D & Среднее \\
\midrule
S. Naive       & 16.42 &  8.31 & 20.40 & 17.39 & 18.70 &  2.80 & 14.00 \\
Auto-ARIMA     & 12.64 &  8.36 & 26.38 & 15.39 & 18.99 &  2.45 & 14.03 \\
ETS            & 12.84 &  7.80 & 21.04 & 16.72 & 15.40 &  2.46 & \textbf{12.71} \\
Prophet        & 14.96 & 17.15 & 24.13 & 22.03 & 20.46 & 10.86 & 18.26 \\
\midrule
DLinear        & 37.64 & 19.32 & 24.89 & 25.70 & 29.51 &  2.43 & 23.25 \\
N-BEATS        & 27.37 & 11.58 & 20.58 & 28.83 & 26.31 &  3.36 & 19.67 \\
TimesNet       & 18.31 &  8.18 & 27.12 & 18.17 & 28.03 &  2.87 & 17.11 \\
\midrule
Autoformer     & 29.31 & 14.63 & 24.66 & 34.17 & 33.86 &  4.02 & 23.44 \\
FEDformer      & 27.75 & 16.91 & 26.64 & 37.00 & 29.94 &  3.43 & 23.61 \\
PatchTST       & 30.02 & 11.29 & 23.72 & 25.00 & 28.87 &  2.85 & 20.29 \\
\bottomrule
\end{tabular}
\end{table}

Средний ранг по~шести категориям (табл.~\ref{tab:mean-rank}) дает
устойчивую ранжировку моделей с~учетом вариативности между наборами данных.

\begin{table}[H]
\centering
\caption{Средний ранг моделей по~шести категориям (sMAPE, меньше --- лучше)}
\label{tab:mean-rank}
\small
\begin{tabular}{@{}lc@{}}
\toprule
Модель & Средний ранг \\
\midrule
ETS            & \textbf{2.00} \\
Seasonal Naive & 2.83 \\
Auto-ARIMA     & 3.17 \\
TimesNet       & 5.50 \\
N-BEATS        & 5.67 \\
Prophet        & 6.00 \\
PatchTST       & 6.00 \\
DLinear        & 7.17 \\
Autoformer     & 8.17 \\
FEDformer      & 8.50 \\
\bottomrule
\end{tabular}
\end{table}

\begin{figure}[H]
\centering
\includegraphics[width=0.9\textwidth]{heatmap_smape.png}
\caption{Тепловая карта sMAPE: модели $\times$ категории}
\label{fig:heatmap-smape}
\end{figure}

\begin{figure}[H]
\centering
\includegraphics[width=0.9\textwidth]{heatmap_rank.png}
\caption{Тепловая карта рангов: модели $\times$ категории}
\label{fig:heatmap-rank}
\end{figure}

Анализ табл.~\ref{tab:cross-smape} и~\ref{tab:mean-rank} позволяет
выделить следующие закономерности.

\textit{Длина ряда как ключевой фактор.}
Классические модели (ETS, Auto-ARIMA) доминируют на~коротких рядах
(yearly, quarterly), но~их превосходство уменьшается по~мере
увеличения длины: на~M3 monthly N-BEATS уже опережает ETS, а~на
M4 daily DLinear лидирует. Это наблюдение согласуется с~тезисом
Зенга и~соавторов~\cite{zeng2023transformers} о~том, что простые
архитектуры могут быть предпочтительнее при достаточных данных.

\textit{Тип декомпозиции и~частота.}
На~квартальных рядах обоих наборов TimesNet (обучаемая 2D-декомпозиция)
показывает результаты на~уровне классических моделей
(8.18~vs.~7.80 на~M3-Q; 18.17~vs.~16.72 на~M4-Q). Это свидетельствует
о~том, что обучаемая декомпозиция становится эффективной при наличии
выраженной сезонной структуры и~хотя бы~40--60~наблюдений.

\textit{Трансформерные архитектуры.}
Autoformer и~FEDformer стабильно занимают нижние позиции
(ранги 8--9 из~10), что объясняется двумя факторами: (1)~обучаемая
декомпозиция серийного типа требует значительного объема данных,
которого нет при per-series обучении на~коротких рядах; (2)~механизм
внимания с~автокорреляцией (Autoformer) или частотным преобразованием
(FEDformer) вносит избыточную сложность для одномерного прогнозирования.
PatchTST, отказавшийся от~явной декомпозиции в~пользу патчей,
показывает заметно лучший результат (ранг~7, sMAPE 20.29) среди
трансформеров.

% ---------------------------------------------------------------------------
\subsection{Helformer: гибридная фиксированная декомпозиция}
\label{subsec:helformer}

Helformer~\cite{kehinde2025helformer} представляет альтернативный подход
к~интеграции декомпозиции в~нейросетевую архитектуру: вместо обучаемого
выделения компонент используется фиксированная декомпозиция
Хольта--Уинтерса (HW). Остаток после снятия тренда и~сезонности
обрабатывается трансформерным энкодером, а~итоговый прогноз формируется
как сумма прогноза HW и~нейросетевой коррекции.

Гибридный характер модели позволяет ожидать лучшей устойчивости
на~коротких рядах: статистическая компонента (HW) обеспечивает
разумный базовый прогноз, а~нейросетевая компонента моделирует
нелинейные остатки. Результаты приведены в~табл.~\ref{tab:helformer-m3}.

\begin{table}[H]
\centering
\caption{Helformer на~M3: sMAPE (\%) в~сравнении с~лучшими моделями}
\label{tab:helformer-m3}
\small
\begin{tabular}{@{}lcccc@{}}
\toprule
Модель & Yearly & Quarterly & Monthly & Overall \\
\midrule
ETS            & 12.84 & 7.80 & 21.04 & 13.89 \\
Auto-ARIMA     & 12.64 & 8.36 & 26.38 & 15.79 \\
TimesNet       & 18.31 & 8.18 & 27.12 & 17.87 \\
N-BEATS        & 27.37 & 11.58 & 20.58 & 19.84 \\
\midrule
Helformer      & \textbf{8.61}$^*$ & \textbf{7.10}$^\dagger$ & \textbf{16.12} & \textbf{11.68}$^\ddagger$ \\
\bottomrule
\multicolumn{5}{l}{\footnotesize $^*$~только 3 ряда; $^\dagger$~27 из 30 рядов; $^\ddagger$~60 из 92 рядов}
\end{tabular}
\end{table}

На~M3 Helformer показывает существенно более низкий sMAPE: 8.61 на~yearly
(против 12.64 у~Auto-ARIMA), 7.10 на~quarterly (против 7.80 у~ETS)
и~16.12 на~monthly (против 20.58 у~N-BEATS). Общий sMAPE составляет
11.68 --- на~2.21~п.п. ниже, чем у~ETS. Необходимо, однако, отметить
существенное ограничение: на~yearly Helformer дал валидные прогнозы
лишь для 3~из~30 рядов из-за ошибок декомпозиции HW на~слишком коротких
рядах (менее 2~полных сезонных циклов). На~quarterly сходимость составила
27~из~30. Таким образом, результат на~yearly следует интерпретировать
с~осторожностью.

\begin{table}[H]
\centering
\caption{Helformer на~M4: sMAPE (\%) в~сравнении с~основными моделями}
\label{tab:helformer-m4}
\small
\begin{tabular}{@{}lcccc@{}}
\toprule
Модель & Quarterly & Monthly & Daily & Overall \\
\midrule
ETS            & \textbf{16.72} & \textbf{15.40} & \textbf{2.46} & \textbf{11.53} \\
TimesNet       & 18.17 & 28.03 & 2.87 & 16.36 \\
\midrule
Helformer      & 17.38 & 18.09 & 2.67 & 12.71 \\
\bottomrule
\end{tabular}
\end{table}

На~M4 Helformer (12.71 overall) располагается между ETS (11.53) и~TimesNet
(16.36), проигрывая ETS на~quarterly (17.38~vs.~16.72) и~monthly
(18.09~vs.~15.40), но~опережая TimesNet на~monthly на~9.94~п.п.
На~ежедневных рядах Helformer (2.67) практически сопоставим с~ETS (2.46).

\textit{Интерпретация.} Фиксированная декомпозиция HW обеспечивает
стабильный базовый прогноз, который нейросетевая компонента дополняет,
но~не~разрушает. Это объясняет преимущество Helformer на~M3,
где ряды короче: статистическая компонента компенсирует недостаток
данных для обучения нейросети. На~M4, где ряды длиннее, ETS
самостоятельно справляется не~хуже, и~дополнительная сложность
нейросетевой коррекции не~окупается на~всех категориях.

% ---------------------------------------------------------------------------
\subsection{Итеративное прогнозирование}
\label{subsec:rollout}

Прямой прогноз (direct) предполагает генерацию всего горизонта за~один
проход. В~реальных приложениях часто применяется итеративный прогноз
(rollout), при котором модель пошагово продвигается вперед,
используя собственные предсказания в~качестве входных данных.
Итеративный режим подвержен накоплению ошибок: неточность каждого
шага вносит систематическое смещение в~последующие
шаги~\cite{cerqueira2020evaluating}.

Эксперимент по~сравнению direct и~rollout проведен на~обоих
наборах данных (M3: 92~ряда, M4: 90~рядов) для девяти моделей.
Результаты на~M3 представлены в~табл.~\ref{tab:rollout-m3}.

\begin{table}[H]
\centering
\caption{Прямой vs итеративный прогноз на~M3: sMAPE (\%)}
\label{tab:rollout-m3}
\small
\begin{tabular}{@{}lccc@{}}
\toprule
Модель & Direct & Rollout & $\Delta$ \\
\midrule
Seasonal Naive & 14.94 & 14.94 & 0.00 \\
Auto-ARIMA     & 15.79 & 15.79 & 0.00 \\
ETS            & 13.89 & 13.61 & $-$0.28 \\
\midrule
DLinear        & 29.76 & 30.16 & +0.40 \\
N-BEATS        & 18.55 & 22.42 & +3.87 \\
TimesNet       & 17.47 & 21.75 & +4.28 \\
\midrule
Autoformer     & 23.08 & 28.03 & +4.94 \\
FEDformer      & 22.34 & 28.62 & \textbf{+6.28} \\
PatchTST       & 21.33 & 23.16 & +1.84 \\
\bottomrule
\end{tabular}
\end{table}

Классические модели практически нечувствительны к~переходу на~rollout:
Seasonal Naive и~Auto-ARIMA имеют нулевой прирост (их прогнозы
по~определению строятся итеративно), а~ETS даже незначительно
улучшается ($-$0.28). Среди нейросетевых моделей наибольшая
деградация наблюдается у~FEDformer (+6.28~п.п.), за~ним следуют
Autoformer (+4.94) и~TimesNet (+4.28). PatchTST демонстрирует
наименьший прирост ошибки среди нейросетевых моделей (+1.84),
что делает его наиболее устойчивым трансформером для итеративного
применения.

Аналогичная закономерность наблюдается на~M4 (табл.~\ref{tab:rollout-m4}).

\begin{table}[H]
\centering
\caption{Прямой vs итеративный прогноз на~M4: sMAPE (\%)}
\label{tab:rollout-m4}
\small
\begin{tabular}{@{}lccc@{}}
\toprule
Модель & Direct & Rollout & $\Delta$ \\
\midrule
Seasonal Naive & 13.11 & 13.11 & 0.00 \\
Auto-ARIMA     & 12.34 & 12.34 & 0.00 \\
ETS            & 11.58 & 12.41 & +0.83 \\
\midrule
DLinear        & 20.96 & 20.45 & $-$0.51 \\
N-BEATS        & 18.22 & 22.06 & +3.85 \\
TimesNet       & 14.79 & 20.83 & +6.04 \\
\midrule
Autoformer     & 21.97 & 26.27 & +4.30 \\
FEDformer      & 24.01 & 27.75 & +3.74 \\
PatchTST       & 18.33 & 20.84 & +2.51 \\
\bottomrule
\end{tabular}
\end{table}

На~M4 TimesNet демонстрирует наибольшую деградацию (+6.04~п.п.),
превышая FEDformer (+3.74). Это объясняется тем, что
двумерная декомпозиция TimesNet, эффективная при прямом прогнозе,
порождает артефакты при последовательном применении к~зашумленным
входам. PatchTST (+2.51) вновь наиболее устойчив среди трансформеров.
DLinear ($-$0.51) практически нечувствителен к~режиму прогнозирования
благодаря своей линейной архитектуре.

\textit{Обобщение.} Средняя деградация при переходе на~rollout
по~двум наборам данных составляет: FEDformer~$+5.01$, TimesNet~$+5.16$,
Autoformer~$+4.62$, N-BEATS~$+3.86$, PatchTST~$+2.18$, DLinear~$-0.05$.
Модели с~более сложной обучаемой декомпозицией (FEDformer, TimesNet)
накапливают ошибки быстрее, поскольку каждый итеративный шаг проходит
через нелинейное преобразование разложения, усиливающее отклонения.
PatchTST, работающий непосредственно с~патчами входного ряда без
явной декомпозиции, оказывается наиболее робастным.

% ---------------------------------------------------------------------------
\subsection{Предобучение трансформерных моделей}
\label{subsec:pretrain}

Одной из~ключевых гипотез настоящей работы является предположение,
что слабые результаты трансформерных моделей в~per-series режиме
обусловлены не~архитектурными ограничениями, а~недостатком обучающих
данных. Для проверки этой гипотезы проведен эксперимент с~предобучением
(pretrain-finetune) на~месячных рядах M3.

\textbf{Протокол.} На~первом этапе каждая из~трех трансформерных моделей
(PatchTST, Autoformer, FEDformer) предварительно обучается на~обучающих
выборках всех 1428~месячных рядов M3 в~режиме скользящего окна
(порядка 70\,000 обучающих окон, 40~эпох). На~втором этапе
предобученная модель дообучается на~каждом из~30~тестовых рядов
(10~эпох). Результат сравнивается с~обучением <<с~нуля>> (from-scratch,
80~эпох) на~тех же рядах.

\begin{table}[H]
\centering
\caption{Предобучение на~M3 monthly: sMAPE (\%) pretrain vs from-scratch}
\label{tab:pretrain}
\small
\begin{tabular}{@{}lccc@{}}
\toprule
Модель & From-scratch & Pretrain+finetune & Улучшение \\
\midrule
PatchTST   & 19.03 & \textbf{16.99} & $-$10.8\,\% \\
Autoformer & 16.67 & \textbf{15.08} & $-$9.6\,\% \\
FEDformer  & 18.56 & \textbf{17.22} & $-$7.3\,\% \\
\bottomrule
\end{tabular}
\end{table}

Все три модели демонстрируют значимое улучшение при предобучении:
PatchTST снижает sMAPE на~10.8\,\%, Autoformer --- на~9.6\,\%,
FEDformer --- на~7.3\,\%. При этом примечательно, что Autoformer
с~предобучением (15.08) опережает Autoformer from-scratch (16.67)
и~приближается к~уровню ETS (21.04 на~тех же рядах с~учетом
того, что pretrain-выборка включает 32~ряда, частично отличных
от~основных 30).

\textit{Интерпретация.} Результаты подтверждают гипотезу о~дефиците
данных: трансформерные архитектуры способны эффективно использовать
обучаемую декомпозицию и~механизм внимания, но~для этого им
требуется значительно больше данных, чем предоставляет один
ряд~\cite{bandara2020forecasting}. Предобучение на~корпусе рядов
позволяет модели сформировать обобщенные представления о~трендовых
и~сезонных паттернах, которые затем эффективно адаптируются
при дообучении на~конкретном ряде.

Масштаб улучшения (7--11\,\%) указывает на~то, что при полноценном
глобальном обучении (или cross-domain transfer learning)
трансформерные модели потенциально могут существенно сократить
разрыв с~классическими подходами --- и~даже превзойти их,
как продемонстрировано в~ряде недавних
работ~\cite{nie2023patchtst, liu2024itransformer}.

% ---------------------------------------------------------------------------
\subsection{Обсуждение результатов}
\label{subsec:discussion}

Полученные результаты позволяют ответить на~исследовательские
вопросы, сформулированные во~введении.

\textbf{RQ1: Когда обучаемая декомпозиция эффективнее фиксированной?}
Фиксированная декомпозиция (ETS, Helformer/HW) доминирует на~коротких
рядах (менее 50~наблюдений) и~при per-series обучении. Обучаемая
декомпозиция (TimesNet, N-BEATS) становится конкурентоспособной при
длине ряда свыше 80--100~наблюдений (M3 monthly, M4 daily).
Обучаемая декомпозиция трансформерного типа (Autoformer, FEDformer)
показывает преимущества только при наличии предобучения или
глобального обучения.

\textbf{RQ2: Какой механизм учета долгосрочных зависимостей наиболее
эффективен для коротких рядов?}
Для коротких одномерных рядов в~per-series режиме экспоненциальное
сглаживание (ETS) и~классические авторегрессионные модели (ARIMA)
остаются наиболее эффективными. Среди нейросетевых архитектур
TimesNet с~2D-свертками и~N-BEATS с~residual-стеками показывают
лучшие результаты, превосходя трансформерные модели.

\textbf{RQ3: Оправдана ли сложность трансформерных архитектур?}
В~per-series режиме --- нет: PatchTST (ранг~7), Autoformer (ранг~8.17),
FEDformer (ранг~8.50) проигрывают более простым N-BEATS (ранг~5.67)
и~TimesNet (ранг~5.50). Однако эксперимент с~предобучением показывает,
что при доступе к~корпусу данных трансформеры реализуют свой потенциал:
улучшение на~7--11\,\% при предобучении свидетельствует о~том, что
их сложность оправдана в~режиме глобального обучения.

\textbf{RQ4: Влияет ли тип декомпозиции по-разному для разных частот?}
Да, влияние выражено. На~квартальных рядах TimesNet (обучаемая 2D)
практически не~уступает ETS (фиксированная): 8.18~vs.~7.80 на M3,
18.17~vs.~16.72 на~M4. На~годовых рядах любая обучаемая декомпозиция
проигрывает фиксированной. На~ежедневных рядах с~длинными историями
даже простая фиксированная декомпозиция DLinear (MA-скользящее среднее)
достаточна для лидирующего результата.

\textbf{RQ5: Каковы реальные границы per-series обучения нейросетей?}
Эксперименты показывают критическую границу порядка 50--100~наблюдений.
Ниже этого порога нейросетевые модели систематически уступают
статистическим. Выше --- разрыв сокращается и на~отдельных категориях
(M3 monthly, M4 daily) исчезает. Предобучение на~корпусе рядов
позволяет эффективно сдвинуть эту границу вниз.

% ---------------------------------------------------------------------------
\subsection{Ограничения эксперимента}
\label{subsec:limitations}

Полученные результаты следует интерпретировать с~учетом ряда ограничений.

\begin{enumerate}[nosep]
\item \textit{Объем выборки.} 30~рядов на~категорию составляют 3--5\,\%
от~полных наборов M3/M4. Хотя средние значения метрик достаточно
устойчивы, отдельные ряды с~аномальным поведением могут существенно
влиять на~результаты, особенно в~условиях высокой дисперсии sMAPE.

\item \textit{Отсутствие подбора гиперпараметров.} Все нейросетевые модели
используют фиксированные гиперпараметры без автоматического подбора
(no hyperparameter tuning). Классические модели (Auto-ARIMA, ETS)
включают встроенный автоподбор, что дает им естественное преимущество.
Данное неравенство осознанно: оно отражает типичную практику, когда
нейросетевые модели развертываются с~параметрами <<по~умолчанию>>.

\item \textit{Per-series режим.} Нейросетевые модели, в~особенности
трансформерные, проектируются для глобального обучения на~больших
корпусах~\cite{nie2023patchtst, wu2021autoformer}. Per-series режим
заведомо неоптимален для них и~не~отражает лучшие практики промышленного
применения. Эксперимент с~предобучением частично компенсирует
этот фактор, но~полноценное глобальное обучение остается за~рамками
данной работы.

\item \textit{Helformer: неполные данные.} На~годовых рядах M3
Helformer дал валидные результаты только для 3~из~30~рядов.
Это существенно ограничивает репрезентативность его yearly-метрик
и~указывает на~необходимость дополнительного исследования
устойчивости HW-декомпозиции.

\item \textit{Вычислительные ограничения.} Autoformer и~FEDformer
обучались с~30~эпохами на~ежедневных рядах (вместо 80), что может
быть недостаточно для полной сходимости и~частично объясняет их
слабые результаты на~M4 daily.

\item \textit{Единственный прогон.} Все эксперименты выполнены
с~фиксированным random seed (42) без усреднения по~нескольким
запускам. Стохастичность инициализации нейросетевых моделей
может вносить дополнительную вариативность, не~учтенную
в~представленных результатах.
\end{enumerate}
